\section{Microscopic tests and finite-chain controls}
\label{app:ground-controls}

Zero-temperature spectra test the state dependence of exchange
cancellation and separate the composite response from endpoint effects.
The comparisons below combine microscopic exact diagonalization with
effective-chain calculations at matched local roots.

\subsection{Different powers on the same periodic geometry}

Local diagonalization found simple roots at each of 16 sampled values
of $Q/\pi$ from $0.55$ to $1$, using $q=(Q+4\pi)/5$ for the quartet
and $q=(Q+2\pi)/5$ for the doublet [Fig.~\ref{fig:rank-controls}(a)].
All sampled local gaps were positive, with minimum $0.04263$.
The quartet coefficient $b$ was nonzero throughout this sample,
whereas the doublet quadrupole vanished identically.
These local roots provide candidate interference settings; a nonzero
many-body residual must be established separately at each setting.

The microscopic comparison uses the same $N=3$, $L=15$ periodic
geometry for both local spins. Table~\ref{tab:rank-roots} specifies simple
roots at different physical momenta but the same effective momentum
$Q=2\pi/3$. Both roots belong to the Hamiltonian in
Eq.~\eqref{eq:rank-ham}.
The odd periodic geometry is frustrated, so this comparison tests the
finite-system mechanism rather than an infinite unfrustrated chain.

\begin{table}[!t]
\caption{Local roots and predictions for the microscopic three-cell comparison.
The local gap is evaluated at the root. $C_0=\|k\|^2$ and $\Cmin$ are
effective predictions for the coefficient of $\beta^2$, not fits to the
microscopic weights.}
\label{tab:rank-roots}
\begin{ruledtabular}
\begin{tabular}{lrr}
 & Quartet & Doublet\\
$s$ & $3/2$ & $1/2$\\
$q/\pi$ & $14/15$ & $8/15$\\
$\alpha_*$ & $-0.38103454$ & $1.45352917$\\
$\Dloc$ & $0.10336638$ & $1.10511693$\\
$\kappa_{\mathrm{opt}}$ & $-9.45913710$ & $-0.17229706$\\
$C_0$ & $0.54511625$ & $0.000474008$\\
$\Cmin$ & $0.14202789$ & $<10^{-30}$\\
\end{tabular}
\end{ruledtabular}
\end{table}

We diagonalized the microscopic $S^z=1/2$ sector of dimension 6435 at
$g=0.005,0.01,0.02,0.04$, both at the fixed local root and with its
predicted linear retuning. There are two degenerate ground states in this
sector; their equal average excludes the entire ground eigenspace from the
inelastic weight. All ten quartet or one doublet low-band excited states
in this sector were included, with additional states above the cutoff used
to check band completeness. The largest eigen-residual was below
$3.1\times10^{-13}$.

The retuned quartet weight remained quadratic at weak coupling
[Fig.~\ref{fig:rank-evidence}(a)]. At the smallest $g$, its coefficient
$\Wlow/\beta^2=0.14091678$ differs from the independent prediction
$0.14202789$ by $0.7823\%$. A log--log fit over the three smallest couplings
gave a slope $1.9869$. For the doublet the corresponding slope was $4.0167$
[Fig.~\ref{fig:rank-evidence}(b)]. Its $\Wlow/\beta^4$ varied from
$1.0694\times10^{-4}$ to $1.1290\times10^{-4}$ across all four couplings.
These finite-interval slopes test perturbative orders along the specified
retuning trajectories. The doublet result supports
Eq.~\eqref{eq:rank-doublet}; further retuning may suppress its weight
below fourth order.

The cancellation was also checked at the operator level. In the full
eight-dimensional retained doublet space, the relative norm of
$K_q+\kappa_{\mathrm{opt}}A'D_Q$ was $4.32\times10^{-15}$.
The analogous quartet operator did not vanish in its full 64-dimensional
space. A restart with a different random vector, larger Krylov space, and
tighter tolerance reproduced the two weakest doublet weights to relative
differences below $1.3\times10^{-9}$.

\begin{figure*}[t]
\centering
\includegraphics[width=0.98\textwidth]{rank_compensation.pdf}
\caption{Retuned-weight comparison and validation limits.
(a,b) Absolute microscopic low-band weights on the same three-cell periodic
geometry, at the parameters in Table~\ref{tab:rank-roots}. Dashed curves keep
$\alpha=\alpha_*$; solid curves use $\alpha_*+\kappa_{\mathrm{opt}}\beta$.
The vertical ranges differ, but no independent intensity normalization is
applied. Points are calculated values; segments only connect them.
(c) Leading quartet minimum $\Cmin$ at the exact $q=\pi$ root for even
effective chains. ED denotes exact diagonalization, OBC/PBC denote open/periodic
boundaries. The accepted DMRG point is an effective $N=20$ open chain;
the $N=40$ cross failed the preset energy-convergence criterion and is not
used as converged evidence. No thermodynamic fit is drawn.
(d) Microscopic $L=20$ open-chain comparison of the retuned leading probe
with the line-weight vector and integrated weight. The errors at larger
$g$ reveal failure of the leading theory. All panels use calculated data
without smoothing; logarithmic axes in (a,b,d) change only the display.}
\label{fig:rank-evidence}
\end{figure*}

\subsection{Even periodic chains and the open-end contribution}

The six-cell effective ring uses the same local roots and effective
momentum $Q=2\pi/3$ as Table~\ref{tab:rank-roots}.
The quartet effective minimum was $0.06641159$, while the doublet minimum
was below $4\times10^{-32}$. The quartet's optimum changed to
$\kappa=-5.56848$, reflecting its dependence on the many-body ground state;
the doublet's operator-level optimum remained $-0.17229706$.
Thus the contrast also occurs without odd-ring frustration in the effective
model. This is not a microscopic thirty-spin calculation.
At the same quartet root, the accepted $N=12$ periodic MPS gives
$\Cmin=0.08315075$ and $\kappa_{\mathrm{opt}}=-5.67393$
[Fig.~\ref{fig:rank-controls}(d)]. Increasing its bond dimension from
384 to 512 changes this weight by only $6.15\times10^{-7}$ relatively.
Its scalar contribution is removed to within $1.4\times10^{-17}$ in
the residual norm squared, again identifying a composite rather than
endpoint source. These size-dependent coefficients are not extrapolated.

For periodic effective chains at the exact quartet root in
Eq.~\eqref{eq:rank-exact-root}, the sizes $N=6,8,10$ gave
$\Cmin=0.000443300,0.000649072,0.000751090$, respectively.
However, $N=4$ gave zero to numerical precision
[Fig.~\ref{fig:rank-evidence}(c)]. This exception rules out attributing an
unavoidable nonzero residual to the local spin alone.

Periodic DMRG extends this test to $N=12$, where bond dimension 512 gives
$\Cmin=0.00080811794$. Relative to bond dimension 384, the coefficient
changes by $3.52\times10^{-5}$ and the energy per cell by
$2.97\times10^{-9}$, passing the original convergence thresholds.
The energy variance per cell is $6.73\times10^{-10}$, and the bond-energy
spread is $1.41\times10^{-8}$. These periodic results exclude an endpoint
origin for this finite-chain residual without establishing a thermodynamic
lower bound.
The $N=16$ calculation at bond dimension 768 had a stable weight but an
energy change per cell of $1.304\times10^{-8}$, above the preset
$10^{-8}$ gate. It is marked unaccepted in Fig.~\ref{fig:rank-controls}(c)
and is not used to extend the convergence claim.

To distinguish bulk-link physics from endpoints, we separated $K_q$ into the
one-spin and quadrupole-times-spin pieces of Eq.~\eqref{eq:rank-link} before
projecting each perpendicular to $d$. For all periodic cases, the perpendicular
one-spin norm squared was below $2\times10^{-30}$. The nonzero residuals are
therefore due to the composite source. Open chains do not share this property.
At $N=10$ OBC, the one-spin, composite, and interference contributions to
$\Cmin$ were $0.000286554$, $0.001244067$, and $-0.000513469$.
Their sum is $0.001017151$. Calling this whole open-chain residual a pure
composite response would be incorrect.

Retuning also changed which excitations dominated the remaining effective
weight. For $N=10$ PBC, the lowest-line share decreased from $0.94007$ at the
fixed root to $1.04\times10^{-5}$ at the integrated-weight optimum. The mean
energy of the residual, in units of $\beta$, rose from $0.09456$ to $0.85896$.
These numbers describe an effective finite chain. They do not imply a bulk
continuum, and the integrated optimum is not exactly the value that cancels
the lowest line alone.

An independent MPS calculation reproduced the $N=10$ open-chain ED minimum
to relative error $2.4\times10^{-6}$ at bond dimension 128. Refinement to
bond dimension 256 gave an accepted $N=20$ minimum $0.00100937976$.
The $N=40$ weight appeared stable but its energy change per cell,
$1.0814\times10^{-8}$, remained above the preset $10^{-8}$ threshold.
It is reported as unaccepted, not used to infer a nonzero thermodynamic limit.

\subsection{Microscopic validity at the exact quartet root}

\begin{figure*}[t]
\centering
\includegraphics[width=0.98\textwidth]{rank_controls.pdf}
\caption{Local mechanism and periodic-chain controls.
(a) Simple local roots at 16 common effective momenta, with
$q=(Q+4\pi)/5$ for quartets and $q=(Q+2\pi)/5$ for doublets. Segments connect
calculated roots; this is not a many-body phase diagram.
(b) Exact channel coefficient $r_J$ in Eq.~\eqref{eq:bond-channel}, not a
spectral intensity. Different quartet channels cannot all be compensated by
the same scalar shift at the link-operator level.
(c) Quartet $\Cmin$ at $q=\pi$ on even periodic effective chains.
The $N=12$ DMRG point passed the original convergence tests; the $N=16$
cross did not and is not converged evidence. No bulk fit is drawn.
(d) Effective chains at the same roots and $Q=2\pi/3$ as
Table~\ref{tab:rank-roots}; the quartet optimum is recomputed for each size.
The doublet dashed line is the first-order operator-cancellation result,
with ED dots at $N=3,6$, not a microscopic finite-$\beta$ zero.
Panels (c,d) retain the normalization $1/\sqrt{5N}$ but use different physical
momenta and local roots; they are not one common size series.
All weights are computed without smoothing or reflection averaging.}
\label{fig:rank-controls}
\end{figure*}

Four open pentamers provide a second microscopic test at $q=\pi$, without
the frustration of the three-cell ring. The effective prediction for their
retuned coefficient is $\Cmin=0.00067813935$.
At $g=0.000625$, microscopic diagonalization gave $0.00067904306$, an
integrated-weight error of $0.1333\%$. The relative Euclidean error of the
sorted low-band line-weight vector was $1.2254\%$.
Doubling $g$ to $0.00125$ approximately doubled that line error to
$2.4625\%$; an independent solver restart at the smaller $g$ agreed to
relative line error $3.56\times10^{-9}$.

The agreement is restricted to weak coupling. At $g=0.02$ and $0.04$, the
line-vector errors were $46\%$ and $109\%$, and the integrated-weight errors
were $14\%$ and $50\%$ [Fig.~\ref{fig:rank-evidence}(d)]. These results
prevent extending the leading composite-probe interpretation quantitatively
to the much stronger couplings of the original full $L=40$ spectra.
